\documentclass[10pt,a4paper]{article}
\usepackage[utf8]{inputenc}
\usepackage[T1]{fontenc}
\usepackage[ngerman]{babel}
\usepackage[margin=1.7cm,top=1.5cm,bottom=1.7cm]{geometry}
\usepackage{titlesec}
\titleformat*{\section}{\normalfont\large\bfseries}
\titleformat*{\subsection}{\normalfont\normalsize\bfseries}
\titlespacing*{\section}{0pt}{8pt}{3pt}
\titlespacing*{\subsection}{0pt}{6pt}{2pt}
\usepackage{amsmath,amssymb}
\usepackage{listings}
\usepackage{xcolor}
\usepackage{tikz}
\usepackage{fancyhdr}
\usepackage{parskip}
\usepackage{needspace}

\newcommand{\mcbox}{\raisebox{-1pt}{\fbox{\rule{0pt}{7pt}\rule{7pt}{0pt}}}\hspace{6pt}}
\newcommand{\antwortlinien}[1]{%
  \par\vspace{2pt}%
  \newcount\zeilen \zeilen=#1
  \loop
    \noindent\rule{\linewidth}{0.4pt}\par\vspace{11pt}%
    \advance\zeilen by -1
  \ifnum\zeilen>0 \repeat
}
\lstset{language=Python,basicstyle=\ttfamily\footnotesize,keywordstyle=\bfseries,
  commentstyle=\itshape\color{gray},showstringspaces=false,frame=single,
  framerule=0.3pt,xleftmargin=8pt,framexleftmargin=6pt,aboveskip=6pt,belowskip=6pt}

\pagestyle{fancy}
\fancyhf{}
\fancyhead[L]{\small MDT5/2 -- Session 8: SVM, Kernel-Trick, OC-SVM, SVDD}
\fancyhead[R]{\small Übungsaufgaben zur Klausurvorbereitung}
\fancyfoot[C]{\small Seite \thepage}
\renewcommand{\headrulewidth}{0.4pt}

\begin{document}

\begin{center}
  {\Large\bfseries Übungsaufgaben Session 8}\\[4pt]
  {\large Kapitel 8: SVM, Soft Margin, Kernel-Trick, OC-SVM, SVDD -- Praktikum 09}\\[4pt]
  \small Stil der Übungssammlung MDT5/2 -- generierte Aufgaben, keine Originale
\end{center}

\vspace{4pt}

\section*{Aufgabe 1 -- Multiple Choice mit Begründung}
\emph{Markieren Sie jeweils die einzige richtige von drei Aussagen und begründen Sie Ihre Entscheidung.}

\subsection*{1.1 \ Für die Support Vector Machine gilt:}
\mcbox Die optimale Trennebene hängt von allen Trainingspunkten gleichermaßen ab.

\mcbox Die optimale Trennebene maximiert den Abstand zu den nächsten Punkten beider
Klassen und hängt nur von den Support Vektoren ab.

\mcbox Die SVM liefert je nach Initialisierung unterschiedliche Trennebenen.

Begründung:
\antwortlinien{2}

\subsection*{1.2 \ Für den Straffaktor $C$ (Soft Margin) gilt:}
\mcbox Je größer $C$, desto stärker werden Verletzungen der Margin-Bedingungen bestraft,
desto enger folgt die Ebene den Trainingsdaten.

\mcbox Je größer $C$, desto mehr Fehlklassifikationen werden im Training toleriert.

\mcbox $C$ bestimmt die Kernelgröße.

Begründung:
\antwortlinien{2}

\subsection*{1.3 \ Für den RBF-Kernel gilt:}
\mcbox Die zugehörige Transformation bildet alle Punkte auf eine Einheitskugel in einem
Raum mit unendlich vielen Dimensionen ab.

\mcbox Der RBF-Kernel entspricht einer Transformation in einen Raum mit genau drei Dimensionen.

\mcbox Für den RBF-Kernel spielt die Skalierung der Eingabedaten keine Rolle.

Begründung:
\antwortlinien{2}

\subsection*{1.4 \ Für OC-SVM und SVDD gilt:}
\mcbox SVDD findet eine Kugel minimalen Volumens um die Normaldaten; bei Verwendung
des Gauß-Kernels sind SVDD und OC-SVM äquivalent.

\mcbox Die OC-SVM benötigt zwingend gelabelte Anomalien im Training.

\mcbox Die OC-SVM mit linearem Kernel ist in der Praxis die Standardwahl.

Begründung:
\antwortlinien{2}

\section*{Aufgabe 2 -- SVM von Hand (Zwei-Klassen-Fall)}

Sie übernehmen die Rolle einer SVM mit linearem Kernel und sehr hohem Straffaktor
(z.\,B. $C=1000$). Die Trainingspunkte sind unten eingezeichnet
($\times$ = Klasse $+1$, $\bullet$ = Klasse $-1$).

\begin{center}
\begin{tikzpicture}[scale=1.0]
\draw[->] (-0.5,0) -- (6.5,0) node[right] {$x_1$};
\draw[->] (0,-0.5) -- (0,5.5) node[above] {$x_2$};
\draw[help lines,step=1] (0,0) grid (6,5);
\foreach \x in {1,...,6} \node[below] at (\x,0) {\small \x};
\foreach \y in {1,...,5} \node[left] at (0,\y) {\small \y};
\foreach \p in {(3,3),(4,4),(4,2),(5,3),(5,5)} \node at \p {$\times$};
\foreach \p in {(1,1),(0.5,2),(2,0.5),(1.5,1.8),(0.7,0.7)} \fill \p circle (2.5pt);
\end{tikzpicture}
\end{center}

a) Markieren Sie die Support Vektoren, zeichnen Sie die beiden durch sie definierten
Parallelgeraden sowie die optimale Trennungsgerade ein.

b) $C$ wird nun schrittweise stark reduziert. Beschreiben Sie die möglichen Auswirkungen
auf Trenngerade und Klassifikationseigenschaften sowie die Gefahr eines zu kleinen $C$.
\antwortlinien{4}

c) Beschreiben Sie kurz den Kernel-Trick und begründen Sie, warum er in dieser Aufgabe
nicht benötigt wird.
\antwortlinien{3}

d) Ein neuer Punkt liegt bei $(2{,}5,\ 2{,}5)$. Nahe an welcher Seite der Trenngeraden liegt
er, und was folgt daraus für seine Klassifikation?
\antwortlinien{2}

\section*{Aufgabe 3 -- OC-SVM von Hand}

Die 100 normalen Trainingspunkte bilden im Plot eine kompakte Wolke; es gibt keine
zweite Klasse.

\begin{center}
\begin{tikzpicture}[scale=1.0]
\draw[->] (-0.5,0) -- (6.5,0) node[right] {$x_1$};
\draw[->] (0,-0.5) -- (0,5.5) node[above] {$x_2$};
\foreach \p in {(3,3),(3.4,3.2),(2.8,3.5),(3.6,2.8),(3.1,2.6),(2.6,3.1),(3.8,3.4),(3.3,3.8),(2.9,2.9),(3.5,3.6),(2.7,2.7),(3.2,3.3),(3.7,3.1),(3.0,3.6),(2.5,3.3)}
  \fill \p circle (2pt);
\fill (1.2,1.0) circle (2pt);
\node[below right] at (1.2,1.0) {\small $p_1$};
\end{tikzpicture}
\end{center}

a) Skizzieren Sie qualitativ die Entscheidungsgrenze einer OC-SVM mit RBF-Kernel und gut
gewähltem $\nu$ (z.\,B. $\nu = 0{,}05$) und markieren Sie, wo Normaldaten und wo Anomalien
liegen würden.

b) Welche Bedeutung hat $\nu$ (zwei Schranken aus der Vorlesung)? Was passiert qualitativ
bei $\nu = 0{,}5$?
\antwortlinien{4}

c) Der Punkt $p_1$ liegt in den Trainingsdaten. Wie geht die OC-SVM mit solchen
\glqq schwierigen\grqq{} Trainingspunkten um, wenn $\nu$ passend gewählt ist?
\antwortlinien{2}

d) Warum ist die OC-SVM mit \emph{linearem} Kernel meist nicht praktikabel?
(Stichwort Trennung vom Ursprung)
\antwortlinien{3}

e) Warum ist der Kernel-Trick (RBF) für OC-SVMs in der Praxis notwendig?
Gehen Sie auf die Abbildung auf die Einheitskugel ein.
\antwortlinien{4}

\section*{Aufgabe 4 -- Einfluss von $\boldsymbol{\gamma}$}

Der RBF-Kernel lautet $k(\boldsymbol{a},\boldsymbol{b}) = e^{-\gamma\,\|\boldsymbol{a}-\boldsymbol{b}\|^2}$.

a) Was beschreibt $\gamma$ anschaulich (Einfluss der Nachbarschaft)?
\antwortlinien{2}

b) Beschreiben Sie qualitativ die Entscheidungsgrenze einer OC-SVM bei viel zu großem
und bei viel zu kleinem $\gamma$.
\antwortlinien{3}

c) In welchen Intervallen werden $C$ bzw. $\gamma$ bei der Metaparameteroptimierung
typischerweise getestet?
\antwortlinien{2}

\section*{Aufgabe 5 -- Kernel-Trick nachrechnen}

Für die Transformation
$\phi(\boldsymbol{x}) = \bigl(x_1^2,\ x_2^2,\ \sqrt{2}\,x_1 x_2\bigr)^T$
gilt $k(\boldsymbol{a},\boldsymbol{b}) = (\boldsymbol{a}^T\boldsymbol{b})^2 = \phi(\boldsymbol{a})^T\phi(\boldsymbol{b})$.

a) Verifizieren Sie diese Gleichheit für $\boldsymbol{a}=(1,2)^T$ und $\boldsymbol{b}=(3,1)^T$,
indem Sie beide Seiten ausrechnen.
\antwortlinien{3}

b) Worin liegt der rechnerische Vorteil des Kernel-Tricks allgemein?
\antwortlinien{2}

c) An welcher Stelle der SVM-Zielfunktion bzw. Entscheidungsfunktion tauchen die
Trainingspunkte nur noch als Skalarprodukte auf -- und warum ermöglicht genau das den
Kernel-Trick?
\antwortlinien{3}

\section*{Aufgabe 6 -- Eigenschaften einordnen}

Nennen Sie aus der Vorlesung:

a) drei Stärken der SVM,
\antwortlinien{3}

b) drei Schwächen bzw. Vorsichtspunkte (mindestens einen zur Datenvorverarbeitung),
\antwortlinien{3}

c) die passenden scikit-learn-Klassen für: Zwei-Klassen-SVM mit $C$, Zwei-Klassen-SVM mit
$\nu$, One-Class-SVM.
\antwortlinien{2}

\vspace{10pt}
\begin{center}
\footnotesize\itshape
Nach Bearbeitung: Antworten an Claude schicken (Foto genügt) -- Korrektur mit Begründungs-Check.
\end{center}

\end{document}
